% Auto-generated by scripts/06_emit_structure.py
% Source: scans 175–177
% Section id: appendix/1
% Phase 3 deliverable. Footnotes (Phase 4), citations (Phase 5),
% and Wade-Giles → Pinyin (Phase 6) are not yet applied here.

\chapter{Identification of the Inscribed Turtle Shells of Shang}
\label{app:1}
\authorbyline{James F. Berry}

% source: scan 175, printed 157
% intro-echo-dropped: 玉生+|
So ce HOLES

Ne oe

Se
zie
aE UW

BS 2eisw Se ezsrae cra GD 828 =

thirty-six plastrons, Ting. like Bien, established the presence of Ocadia sinensis and Chinemys reevesi

among the An-yang remains. To date, no systematic effort has been made to identify large groups
of remains from the An-vang site.

2. Material and Methods

The identification of turtle skeletons is usually a comparatively simple matter, involving
the comparison of unknown materials with comparable material of known identification and
origin. The identification of the An-yang turtle remains is rendered more difficult by several
factors: (1) relatively little skeletal material of Chinese turtles is available for comparison or
identification; (2) little is known of variation, either geographic or allometric, in Asiatic turtles:
(3) the plastrons and carapaces at An-yang were mostly burned and broken; and (4) most scholars
do not have access to the plastrons themselves, but only to rubbings. Given these difficulties, a
quantitative approach, such as that adopted by Ting Su (1969), is the most promising. Since the
main goal is to identify the turtle species used by the Shang diviners, I shall take as my corpus the
inscribed shells excavated by the Academia Sinica in the An-yang region between 1928 and 1937,
and reproduced in Ping-pien, Yi-pien, and, to a lesser extent, in Chia-pien (on these three collections,
see sec. 5.4). Since this material consists predominantly of plastrons, all identifications will be based
on plastral relationships. ,

The plastron of modern turtles consists of a layer of bone overlaid by a series of epidermal
scutes. The bones and scutes are arranged in species-specific patterns, but the patterns of the
seams between the bones and the seams between the scutes are different (fig. 3). Only relationships
between scute seams (fig. 31: will be considered here because they are the most easily measured.
It is assumed that these scute ratios will vary sufficiently to allow separation of turtle species.
Unfortunately, on many of the Ping-pien and Yi-pien plastrons the seams between the bones and
scutes appear superimposed and cannot be separated (fig. 3); these shells, together with those
which were scraped clean ‘sec. 1.3.2) so that only the bone seams are visible, are excluded from
the present study.

Total plastron length :PL) was measured; this was taken to be the distance from the most
anterior to the most posterior portion of the plastron. The lengths of the following seams were
measured along the midventral line: intergular seam length (IGL) (the seam, which spans the
distance from the most anterior to the most posterior part of the gular scute, is often not apparent) ;
interhumeral seam length ‘IHL); interpectoral seam length (IPL); interabdominal seam length
(IAB) ; interfemoral seam length (IFL); and interanal seam length (IAL); these are the various
seams formed by the joining of the scutes from the right and left sides along the midline. The width
of the plastron was measured in four places: the gular-humeral seam width (PW1), the humeral-pectoral seam width (PW2\_, the abdominal-femoral seam width (PW3), and the femoral-anal seam
width (PW4) (fig. 31).

Two methods were used in the identification of the plastrons. With the first method, a series
of ratios was obtained by dividing each seam length by the total plastron length. This technique
has the effect of standardizing each seam length with respect to the absolute size of each turtle.
Representatives of each turtle species likely to be encountered among the An-yang remains were
measured and ratios formed. The ratio determined to be the most satisfactory for distinguishing
the species was IGL+IHL/PL. Data was then gathered from as many Chia-pien, Ping-pien, and
Yi-pien plastrons as possible, and the IGL+IHL/PL ratios calculated.
% source: scan 176, printed 158
APPENDIX If

The second method considered all seam lengths on one plastron simultaneously by mathematically reducing the total seam length number to a smaller, more workable, number. This
technique, called “‘principal component analysis,’ has been used effectively for discerning groups
among apparently homogeneous clouds of data (Blackith and Reyment [1971]\}.! The comparison
of a number of ratios simultaneously should produce identifications somewhat more reliable than
those obtained by the first method.

From the standpoint of this study, a weakness of this second method is that differences in
size will usually account for most variation. To avoid this bias, the original data matrix was
standardized by using ratios of the seam lengths/plastron lengths of pure data and seam lengths/
plastron length ratios of log-log transformed data. Sokal (1965; has warned that ratios are of
questionable validity for this purpose since there is no reason to assume that the ratios are normally
distributed (principal component analysis assumes normally distributed data), but since so little
is known of variability or allometric relationships, there seems to be no alternative.

The principal component analysis was performed once each on ratios of untreated and log-log
transformed data on the University of Utah Univac 1108 computer using program BMDO1M
(Dixon [1974]). Since the program has no provision for missing data, only essentially complete
plastrons could be used in this analysis. For comparative purposes, ratios from known examples
of Ocadia sinensis, Chinemys reevest, and Mauremys mutica were included with the unknown Chva-pien,
Ping-pien, and Yi-pien data. The values for each turtle on the first principal component were plotted
versus the values on the second component (figs. 32 and 33).

3. Results

The ratios IGL+IHL/PL are presented in table 3 together with a list of identifications
based on similarity of ratios of An-yang material to known material. The validity of the above
identifications is rendered uncertain by the lack of knowledge of variation in Chinese turtles. For
example, the failure of the humeral-pectoral seam to intersect the entoplastral bone has been used
as one of the characters distinguishing Mauremys from Ocadia (for example, Lindholm [1931]).
However, it has been shown by Nakamura (1934) and more recently by Mao (1971) that this
character is variable in both Mauremys and Ocadia. This situation would affect the validity of the
ratio-based identifications in the following way: assuming the position of the entoplastral bone
remains constant, a humeral-pectoral seam which crosses the entoplastral bone will result in a proportionately shorter intergular plus interhumeral seam (smaller IGL+IHL/PL ratio). This would
result in overlap between individual Mauremys and Chinemys or even Ocadia in this character.

1. With the technique of principal component
analysis, the original n x p data matrix (where

= the number of turtles measured) is standardized by centering and normalizing the data matrix.
A matrix R is computed as n — 1 times the correlation matrix. The set of simultaneous equations
(R — Al)v = o [where J = identity matrix]) are
then solved, yielding a set of eigenvalues (A) and
their corresponding eigenvectors (v) which are the
principal components. (For more detailed descriptions of the mathematical procedures, see
Blackith and Reyment [1971]; Cooley and Lohnes
[1971]; Sneath and Sokal [1973]; and Dixon

[1974]).
We can think of the resulting values as being

represented by a cloud of points in n-dimensional
hyperspace. The first principal component is a
line (a vector) oriented in the cloud in such a way
that its values (eigenvalues) account for as much
variation as possible. The second component is
orthogonal (i.e., at right angles) to the first,
oriented to account for as much variation as possible. Each vector will have a value with respect to
a particular case, so that cases (= turtles) can be
ordered with respect to each vector.

了一一一
% source: scan 177, printed 160
Nevertheless, the results of the identifications based on ratios, though of uncertain reliability,
provide firmer identifications than were previously available for a large number of Chia-pien,
Ping-pien, and Yi-pien plastrons. Most of the plastrons are identified by this first technique as Ocadia
sinensis and Chinemys reevesi.

It can be seen from figs. 32 and 33 that the data assessed by the second technique (principal
component analysis) fall generally into two groups: one group which clusters with Ocadia sinensis
and a second group which clusters with Chinemys reevesi. The identifications based on this analysis
are listed in table 4 along with individual values on the first two principal components. Questionable
identifications are noted with a question mark. Though fewer in number, the identifications based
on this analysis agree well with those based on the ratio analysis.

Ping-pien 184 can be identified with certainty as Testudo emys. As noted by Wu (1943), the
failure of the pectoral scutes to meet at the midline has been used as a diagnostic character for this
species (Boulenger [1889]; Bourret [1941]; Smith [1931]; and Wermuth and Mertens [1961]).

4. Discussion

Plastrons of three species of aquatic freshwater turtles and one terrestrial tortoise have
been identified from the shells used by An-yang diviners.

1. Ocadia sinensis, which occurs on Taiwan, mainland China south of Fukien, Hainan, and
northern Indochina (Schmidt [1927]; Pope [1935]; McDowell [1964]).

2. Chinemys (=Geoclemmys) reevesi, occuring in Korea, Japan, and Taiwan; in north and
central China from southeastern Shantung westward to the Wei Valley in Shensi, throughout the
Yangtze River drainage from Szechwan eastward; and along the southern coast of China to
Kwangtung (Stejneger [1907]; Pope [1935]; McDowell [1964]; Mao \{1971]). Pope reports that
Chinemys is commonly sold in Chinese markets, resulting in probable artificial dispersal by man.

3. Mauremys (=Clemmys) mutica which occurs in Japan, southern China, Hainan, and
Vietnam (Stejneger [1907]; Pope [1935]; McDowell [1964]; and Mao [1971]).

4. Testudo emys, a tortoise which occurs in southeast Asia and the Malay Archipelago where
it is prized for its flesh (Smith [1931]; Pope [1935]).

Notes on the habits of these turtles can be found in Smith (1931), Pritchard (1967), and Mao
(1971). Individuals of Ocadia, Chinemys, and Mauremys have been maintained in excellent condition
for a number of years in the live animal colony of the Herpetology Laboratory of the University
of Utah, where they feed readily on lettuce, fruit, and chopped fish.

The identification of the type specimen of Testudo anyangensis Ping remains in doubt. It was
correctly labeled an aquatic turtle and not a land tortoise by Lindholm ([1931]: Pseudocadia
anyangensis) and independently by Auffenberg ([1962]: Clemmys anyangensis). Bien (1937) referred
to the shell as Ocadia sinensis, and McDowell (1964) referred to it as Mauremys mutica. Nakamura
(1934) and Mao (1971) have shown that the characters used by Ping (1930), Lindholm (1931),
and Bien (1937) are subject to extreme geographic variation, but my ratio IGL+IHL/PL (based
on Ping’s figure reproduced in Pope [1935]) is .154, which, based on the results of this study,
suggests that the shell is Ocadia. In any event, this particular specimen was not used by the Shang
diviners.
